%% ch12 --- Observability and operations: structlog, OpenTelemetry, containers and audits
%%
%% Written September 2026. Every claim about a library here was checked
%% against the installed, pinned package; the two that contradicted the
%% brief are recorded in CLAUDE.md under *Resolved questions*.

\chapter{Observability and operations: structlog, OpenTelemetry, containers and audits}
\label{ch:observability}

\noindent
Your operational reflexes are right and they transfer whole. Structured
logs rather than formatted strings; one identifier threaded through
everything a request touches; spans around the work that costs money; an
image small enough to pull quickly; a scan that tells you which dependency
has a published vulnerability; a shutdown that finishes what it started.
Python has all six, and for each one there is a package that a lot of
people use and broadly agree about.

What does not transfer is the defaults. In .NET the defaults for these six
are the ones a server wants, because the host builder was written for
servers. In Python most of them are the ones an interactive session wants,
because the module they live in is older than the deployment model \dash{}
and the gap is almost never an error. A log configuration that silently
does nothing, a correlation identifier that silently belongs to another
request, an image that silently ships your test dependencies: each of
these runs, and returns, and is wrong.

This chapter assumes Chapter~\ref{ch:services}, which owns FastAPI, its
\api{lifespan}, and configuration read by \pkg{pydantic-settings}; and
Chapter~\ref{ch:testing}, which owns pytest. It owns what happens to the
process afterwards.

\section{The configuration that happens once, to whoever is first}
\label{sec:ch12-logging}

Read the next listing before you run it, and decide what the second call
to \api{basicConfig} does \dash{} it asks for a level and a format that
the first call did not.

\pyfile{code/ch12/logging_default.py}{\code{basicConfig} is not \code{ILogger} configuration}{lst:ch12-logging-default}

\noindent
It does nothing:

\transcript{ch12-logging-default}

\noindent
No exception, no warning, no return value to inspect. The requested
\code{DEBUG} line never appears and the warning that does appear is in the
\emph{default} format, not the one asked for. The rule is one line of the
standard library: \api{basicConfig} configures the root logger only when
the root logger has no handlers, and the first call put one there.

\begin{csbox}
\code{ILogger} is configured once, by you, on the host builder, and handed
to everything else by injection. Nothing a package does at load time can
take that decision away from you, because there is no ambient logger to
take it with. Python's \pkg{logging} is module-level state, and every
importer shares it.
\end{csbox}

\noindent
In the listing the first call is explicit so that you can see it. In a
service it is not yours at all: \api{logging.warning(...)} called at import
time, by anything, calls \api{basicConfig()} for you with no arguments.
Whether your configuration takes effect then depends on import order.

\begin{trapbox}
\textbf{\enquote{\code{logging.basicConfig} and I am done.}} It is a
one-shot, and the shot may already have been fired by a dependency you
have never read. \code{force=True} is the standard library's own answer
\dash{} drop every existing handler and configure from scratch \dash{} and
it is why the next section's bootstrap assigns \code{root.handlers}
outright rather than appending to it.
\end{trapbox}

\noindent
There is a second, smaller reason people replace this module: its default
destination is standard error, and its default format is a string. A
server wants standard output and a machine-readable record.

\section{One line, one dict}
\label{sec:ch12-structlog}

\pkg{structlog} keeps a log record as a dictionary until the last possible
moment. Everything between the call and the output is a \emph{processor}
\dash{} a function that takes the dictionary and returns it \dash{} and the
last processor in the chain renders it. That is the whole design, and it is
the reason the bootstrap below is worth copying rather than describing.

\mermaidfig{obs-log-pipeline}{Two sources of records, one chain and one renderer: the \code{foreign\_pre\_chain} is what puts a library's line through the same processors as your own}{the log pipeline}

\pyfile{code/ch12/log_setup.py}{The bootstrap: configure once, render once, and catch the records that never met structlog}{lst:ch12-log-setup}

\noindent
The part that is easy to leave out is \api{foreign\_pre\_chain}. Without
it your own lines are JSON with a request identifier on them and a
library's lines are bare strings; with it, both go through the same
processors and come out the same shape:

\transcript{ch12-log-setup}

\noindent
The middle line is \pkg{httpx}'s. It was written with the standard
library, by code that knows nothing about structlog and nothing about the
request identifier, and it carries the identifier anyway.

\begin{note}
The demo asks for no timestamp, which is why its lines fit this page. It
is also a defensible production choice: under a collector that stamps
every line it receives, the application's own timestamp is a second and
slightly different answer to a question somebody has already answered.
Keep the default when you are writing to a file.
\end{note}

\noindent
One more thing belongs in the chain rather than in a review checklist.
Chapter~\ref{ch:services} loads configuration with \pkg{pydantic-settings};
declare the secret ones as \api{SecretStr} and they cannot reach a log
line. \api{SecretStr} renders as a row of asterisks in \api{repr}, and
structlog's JSON renderer falls back to \api{repr} for anything the
serialiser refuses \dash{} so a secret logged by accident prints
\code{SecretStr('**********')} where a plain \code{str} prints the secret.
The type is the control.

\section{The identifier two requests must not share}
\label{sec:ch12-contextvars}

Here is the next one to read before running. Two handlers run at the same
time; each writes its own identifier to a module-level variable and to a
\api{ContextVar}, waits, and reads both back. Write down what each will
report.

\pyfile{code/ch12/correlation.py}{Where a per-request value can live when two requests share a thread}{lst:ch12-correlation}

\transcript{ch12-correlation}

\noindent
Handler~A wrote \code{A} and read back \code{B}. Every log line A goes on
to write is filed under B's request, and nothing anywhere reports a
problem. The \api{ContextVar} column is right for both, and the caller's
row shows the other half of the behaviour: a task runs in a \emph{copy} of
its parent's context, so a child's value cannot leak back out.

\begin{csbox}
\api{ContextVar} is \code{AsyncLocal<T>}, closely enough that the mental
model transfers without adjustment. What has no equivalent is the habit
that a static field is safe because you only ever have one request in
flight per thread \dash{} here one thread carries all of them.
\end{csbox}

\begin{trapbox}
\textbf{\enquote{Correlation identifier in a static field, like
\code{AsyncLocal}.}} A module-level variable is a static field, and it is
shared by every request on the loop. The failure is silent, it needs
concurrency to appear, and a single-request test will never show it.
\end{trapbox}

\noindent
In practice you rarely write the \api{ContextVar} yourself:
\api{bind\_contextvars} holds one for you, and
\api{merge\_contextvars} \dash{} the first processor in the chain above
\dash{} is what copies it onto every record. Bind the identifier once per
request and every line for the rest of that request carries it. How the
loop schedules those requests is Chapter~\ref{ch:asyncio}'s subject and is
not repeated here.

\section{Spans: three nouns you already have}
\label{sec:ch12-otel}

OpenTelemetry is the same specification in both ecosystems, so this is
mostly a renaming exercise:

\begin{center}
\small
\begin{tabularx}{\linewidth}{@{}lX@{}}
\toprule
\textbf{.NET} & \textbf{Python} \\
\midrule
\code{ActivitySource} & a \api{Tracer}, from \api{trace.get\_tracer(\_\_name\_\_)} \\
\code{Activity} & a span, from \api{tracer.start\_as\_current\_span()} \\
\code{ActivityListener} plus exporters & a \api{TracerProvider} with span processors \\
\code{Activity.Current} & \api{trace.get\_current\_span()} \\
Resource attributes on the builder & a \api{Resource} passed to the provider \\
\bottomrule
\end{tabularx}
\end{center}

\noindent
The one default worth knowing before you ship is the resource. Create a
provider without one and the SDK fills in \code{service.name} for you. The
value it fills in is the literal string \code{unknown\_service}, so every
span your service has ever emitted arrives in the backend under a name that
is also every other unconfigured service's name.

\pyfile{code/ch12/tracing.py}{A provider, a named service, and the processor that puts the trace identifier on every log line}{lst:ch12-tracing}

\noindent
\api{add\_trace\_ids} is the part to copy. Put it in the chain from
Listing~\ref{lst:ch12-log-setup} and a log line written inside a span
carries the identifiers the trace was recorded under, which is what makes
a line in a log search and a span in a trace viewer one click apart. The
\api{is\_valid} guard matters: outside a span the SDK hands back a context
whose identifiers are zero, and a log field reading all zeros is worse
than an absent one, because it looks like a value.

\begin{versionbox}
Automatic instrumentation \dash{} the \code{opentelemetry-instrument}
wrapper that patches a framework's calls without a code change \dash{} is
real and is the usual starting point, but it lives in distributions this
book does not pin, and which of them covers which library moves faster
than a book can. Nothing here measures it. The manual span above is what
the SDK alone gives you.
\end{versionbox}

\section{The image}
\label{sec:ch12-image}

Three shapes, and the difference between the first and the last is not
subtle. The first is what the reflex writes.

\pyfile{code/ch12/Dockerfile.naive}{Shape one: the full base image, \code{pip}, no lockfile, root}{lst:ch12-dockerfile-naive}

\pyfile{code/ch12/Dockerfile.slim}{Shape two: slim, uv, one stage \dash{} still shipping the tool that did the installing}{lst:ch12-dockerfile-slim}

\pyfile{code/ch12/Dockerfile}{Shape three: two stages, and the one to copy}{lst:ch12-dockerfile}

\begin{verifybox}
These three were not built. The machine this book is compiled on has a
Docker daemon and no route to any registry \dash{} both Docker Hub's and
GitHub's container blob stores answer \code{403} to its egress proxy
\dash{} so no base image can be fetched. Every flag and every environment
variable in them was checked against the installed \code{uv} and the
pinned interpreter; the files as a whole were not executed, and experiment
E7 below is the thing that would execute them.
\end{verifybox}

\mermaidfig{obs-image-stages}{What crosses between the stages is the virtual environment and nothing else: no package manager, no lockfile, no build cache}{the two stages}

\noindent
Four decisions in shape three are worth naming.

\textbf{Two stages.} The builder needs \code{uv} and the lockfile; the
runtime needs neither. Copying only the virtual environment across leaves
the shipped image with no package manager in it, which is a smaller attack
surface as much as it is a smaller image.

\textbf{\code{-{}-frozen} against \code{-{}-locked}.} They are not
synonyms. \code{-{}-locked} \emph{asserts} that \code{uv.lock} is
consistent with \code{pyproject.toml} and fails if it is not;
\code{-{}-frozen} installs what the lockfile says without checking.
Assert in CI, where a stale lock is a defect you want to hear about;
install frozen in the image, where you have already heard.

\code{UV\_COMPILE\_BYTECODE} and \code{PYTHONDONTWRITEBYTECODE} point
opposite ways on purpose. The first, in the builder, compiles every module
once at build time into an immutable layer, so start-up does not pay for
it. The second, in the runtime, stops the process writing any more, which
would only dirty the container's writable layer. Both are right; they are
about different moments.

\textbf{\code{PYTHONUNBUFFERED=1}.} Standard output is block-buffered when
it is a pipe, and in a container it is always a pipe. Without this your
log lines sit in a buffer, arrive in bursts, and are lost outright if the
process is killed \dash{} which is precisely the process you most wanted
the logs from.

\begin{trapbox}
\textbf{\enquote{\code{FROM python:\pyver} and \code{pip install} in the
Dockerfile.}} It works, and it ships a compiler toolchain, a package
manager, your development dependencies and a root user, resolved against
whatever the index served that morning rather than against a lockfile.
Multi-stage, \code{uv sync -{}-frozen -{}-no-dev}, and a numeric
non-root user.
\end{trapbox}

\subsection{What \code{-{}-no-dev} is worth, and what has not been measured}
\label{sec:ch12-measured}

One of those four decisions can be priced without building anything.
\code{uv.lock} records every wheel's size, so the runtime dependency set
and the full one can both be totalled by reading a committed file:

\begin{center}
\small
\begin{tabularx}{\linewidth}{@{}Xrr@{}}
\toprule
& \textbf{Packages} & \textbf{Wheels, MiB} \\
\midrule
Runtime only, as \code{-{}-no-dev} installs & \val{ch12.deps.runtime.count} & \val{ch12.deps.runtime.mib} \\
Everything, as a plain \code{uv sync} installs & \val{ch12.deps.full.count} & \val{ch12.deps.full.mib} \\
\midrule
What the flag removes & \val{ch12.deps.dev.count} & \val{ch12.deps.dev.mib} \\
\bottomrule
\end{tabularx}
\end{center}

\noindent
Leaving the flag out multiplies the dependency payload by
\val{ch12.deps.ratio}, for this book's own lockfile. Two honest
qualifications: those are \emph{download} bytes, and a wheel unpacks to
more than it weighs, so the footprint removed from an image is larger than
the figure rather than smaller; and it is one of three differences between
the shapes, not the image size.

The image size is experiment~E7 \dash{} three shapes built, measured and
started \dash{} and \textbf{it has not been run}, for the reason the box
above gives. Its script is written and committed; it needs a machine with
a container registry in reach. Until it runs, every comparison this
chapter makes between the three shapes is an argument from what they
contain, which is judgement, and there is no table of sizes here because
there is nothing to put in it.

\section{Auditing the lockfile rather than the machine}
\label{sec:ch12-audit}

\code{dotnet list package -{}-vulnerable} reads your project. \pkg{pip-audit}
by default reads the \emph{environment} it is run in, which is a different
question with a similar answer, and one you cannot ask until you have
installed something.

It does have a lockfile mode, and the name is a trap: \code{-{}-locked}
audits a PEP~751 \code{pylock.toml}, and \code{uv.lock} is not one
\dash{} pointed at a uv project it reports that it found no lockfile at
all. The bridge is one command, and then the two halves fit:

\begin{shellcmd}
uv export --no-dev --format pylock.toml -o pylock.toml
uv run pip-audit --locked .
\end{shellcmd}

\noindent
Worth knowing before you put it in CI: your own package is skipped,
because it is not on an index to be looked up, and a skip is not a
failure. \code{-{}-strict} makes it one. Audit the \emph{runtime} set,
as above, or you will be triaging advisories against a linter you do not
ship.

\section{The signal, the drain and the two probes}
\label{sec:ch12-shutdown}

The sequence uvicorn runs when it receives \code{SIGTERM} is in its own
source, and it is worth reading there rather than trusting a summary:
stop accepting new connections, ask the open ones to finish, wait for the
in-flight tasks, and only then run the \api{lifespan}'s shutdown half. So
the cleanup Chapter~\ref{ch:services} put in the lifespan runs after the
requests have drained, which is the order you want.

\begin{warning}
Uvicorn's graceful-shutdown timeout is unset by default, and unset means
\emph{wait indefinitely}. One request that never finishes and the process
never exits, until the orchestrator loses patience and sends
\code{SIGKILL} \dash{} at which point nothing in the lifespan has run at
all. Set \code{timeout\_graceful\_shutdown} to a number smaller than your
orchestrator's grace period.
\end{warning}

\noindent
What that sequence does not do is tell the load balancer. Between the
signal arriving and the balancer noticing, every request it still routes
to this process meets a closed socket. So readiness has to fail first, and
the drain waits out the balancer's polling interval before it starts.

\mermaidfig{obs-shutdown}{Readiness goes first and liveness stays true: a draining process is leaving, not sick, and a failed liveness probe would have it killed mid-drain}{the shutdown sequence}

\pyfile{code/ch12/shutdown.py}{The order, made explicit: readiness withdrawn, then the drain, then the lifespan}{lst:ch12-shutdown}

\transcript{ch12-shutdown}

\noindent
The two probes answer different questions and an orchestrator treats them
completely differently. Liveness failing gets the process killed and
replaced; readiness failing gets it taken out of rotation and left alone.
During a shutdown you want exactly one of those, which is why
\code{live} stays true for the whole of the drain above.

\section{What you have now}
\label{sec:ch12-summary}

Six reflexes, six names. Structured logging is \pkg{structlog} with one
renderer and a \api{foreign\_pre\_chain} so that other people's lines
obey it too. The correlation identifier is a \api{ContextVar}, which is
\code{AsyncLocal} under a different name and is the one place a
module-level variable will cost you a production incident. Tracing is the
same specification you already know, with a resource you must remember to
set. The image is two stages, a lockfile and a user id. The audit is
\pkg{pip-audit} over an exported \code{pylock.toml} rather than over a
machine. And the shutdown is a signal, a withdrawn readiness probe and a
drain, in that order.

Every default in that list was against you, and not one of them would have
raised an error.

\begin{exercise}{e12_01_bootstrap}{One renderer for your lines and everybody else's}
Finish \api{configure\_logging} so that a record written through structlog
and a record written through the standard library both reach the same
stream, both as JSON, and both carrying whatever was bound with
\api{bind\_contextvars}. The third test is the one that catches a
half-finished answer: it is easy to get your own lines right and leave a
library's with no context at all. Write it before you read
Listing~\ref{lst:ch12-log-setup}.
\end{exercise}

\begin{exercise}{e12_02_correlation}{A request identifier two requests cannot share}
\code{RequestContext} is a static field carried across from C\# without
translation, and under any concurrency at all the identifier a handler
reads back is whichever handler wrote last. Keep the three method names
and make the value per-context. The test runs two handlers at once and
also checks the half people forget: a value set inside a task must not
leak back out to the caller.
\end{exercise}

\begin{exercise}{e12_03_trace_id}{The processor that joins a log line to a trace}
Write the processor that stamps the current span's identifiers onto a
record. Inside a span both must appear, as hexadecimal of the right width;
outside one, neither may. The second half is the point: a span context
that is not valid still has identifiers, they are all zeros, and a log
field full of zeros is worse than a missing one.
\end{exercise}

\begin{exercise}{e12_04_readiness}{The order a shutdown has to happen in}
Finish \code{Lifecycle} so that beginning a shutdown withdraws readiness
and leaves liveness alone, and so that the process may exit only once the
shutdown has begun and nothing is in flight. Getting the liveness half
wrong is the interesting failure: it gets a draining process killed with
requests still on it, which is the outage the drain existed to prevent.
\end{exercise}
